\section{Introduction}

LLM-based agents are increasingly augmented with episodic memory: systems that store past experiences and retrieve them to guide future decisions.
Recent work has demonstrated impressive gains from such augmentation (Voyager~\citep{wang2023voyager}, GITM~\citep{zhu2023gitm}, MemoryBank~\citep{zhong2024memorybank}, and Generative Agents~\citep{park2023generative}) under the implicit assumption that storing and retrieving past experience will improve performance.
But does it always?

We investigate this question through a controlled comparative study across two interactive environments: ALFWorld~\citep{shridhar2021alfworld}, a structured household task environment with compositional goals and a fixed action grammar, and WebShop~\citep{yao2022webshop}, an open-ended e-commerce environment with free-form navigation.
In our experiments, the answer depended critically on environment properties.
On ALFWorld, our memory system improves over the ReAct~\citep{yao2023react} baseline by up to $+19.4$pp on unseen tasks, nearly matching Reflexion~\citep{shinn2023reflexion} ($+21.6$pp) at roughly one-third the step cost.
On WebShop, the \emph{same} memory system \emph{degrades} performance by up to $30.5$pp.
Strikingly, a hard-negative control that retrieves maximally \emph{irrelevant} experiences ($-12.5$pp) outperforms the properly configured system ($-30.5$pp), revealing that semantically similar but non-transferable memories are more harmful than obviously irrelevant ones in this setting.
This ${\sim}47$pp swing, consistent across development and held-out splits, is the central empirical finding of this paper.

Our analysis attributes this divergence to three environment properties: \emph{task compositionality} (shared sub-goals vs.\ unique tasks), \emph{action space regularity} (fixed grammar vs.\ free-form), and \emph{learning transfer radius} (one learning benefiting many tasks vs.\ near-zero reuse). In the two environments we study, when all three favored transfer, memory yielded substantial gains; when they did not, retrieved experiences actively misled the agent. Our results suggest that high-similarity retrieval alone does not guarantee performance gains, as retrieved trajectories may lack transferability to the current task — though the generality of this pattern awaits validation in further settings.
This points to a potentially important distinction between retrieval relevance and actionable utility, one that deserves greater attention in future work on memory-augmented agents.
Accordingly, the novelty of this work is primarily empirical and diagnostic: rather than proposing a fundamentally new memory architecture, we use a controlled episodic-memory setup to study when retrieval transfer helps and when it fails.

In summary, our contributions are:
\begin{enumerate}
    \item A \textbf{controlled empirical study} showing opposite effects of the same memory system: $+19.4$pp on ALFWorld vs.\ $-30.5$pp on WebShop, with hard-negative ablations revealing that ``relevant'' but non-transferable memories cause more harm than irrelevant ones.
    \item \textbf{A diagnostic framework for retrieval transfer}, based on three environment properties: task compositionality, action space regularity, and learning transfer radius, that characterize when episodic memory is likely to help or hurt LLM agents.
    \item A \textbf{memory-centric agent framework} with outcome-validated episodic memory and two retrieval mechanisms (proactive CR and on-demand TR), evaluated across two diverse environments.
\end{enumerate}
